\section{二项式定理}

\subsection{二项式定理的定义}

\begin{definition}
    二项式定理是指对于任意整数 $n \geq 0$，有
    \[
        (a + b)^n = \sum_{k=0}^{n} C_n^k a^{n-k} b^k
    \]
    其中 $C_n^k$ 是二项式系数，表示从 $n$ 个不同元素中选取 $k$ 个元素的组合数。
    二项式系数 $C_n^k$ 的计算公式为
    \[
        C_n^k = \frac{n!}{k!(n-k)!}
    \]
    二项式定理可以用于展开二项式的幂次，计算多项式的系数等。
\end{definition}

这里给出一道例题，只需要记住，二项式展开后的第 $r+1$ 项为 $T_{r+1} = C_n^r a^{n-r} b^r$，所以如果题目问第 5 项是多少，那么公式中 $r$ 就是 $4$。

\begin{example}
    在 $\left(x+\frac{2}{x^{2}}\right)^{6}$ 的展开式中，常数项为 \underline{\hspace{2cm}}. (用数字作答)
\end{example}

\begin{solution}
    由 $\left(x+\frac{2}{x^{2}}\right)^{6}$ 的展开式的通项为
    $$T_{k+1} = \mathrm{C}_{6}^{k} x^{6-k} \left(\frac{2}{x^{2}}\right)^{k} = \mathrm{C}_{6}^{k} 2^{k} x^{6-k} x^{-2k} = \mathrm{C}_{6}^{k} 2^{k} x^{6-3k}$$
    令 $6-3k=0$，解得 $k=2$。
    则常数项为 $T_{2+1} = T_{3} = \mathrm{C}_{6}^{2} 2^{2} x^{6-3(2)} = \mathrm{C}_{6}^{2} \cdot 4 \cdot x^{0}$。
    因为 $\mathrm{C}_{6}^{2} = \frac{6 \times 5}{2 \times 1} = 15$，
    所以 $T_{3} = 15 \times 4 = 60$。
    即在 $\left(x+\frac{2}{x^{2}}\right)^{6}$ 的展开式中，常数项为60。
\end{solution}

\vspace{0.8cm}

\begin{problem}
$(a+b)^n$ 展开式中共有 $n$ 项，是否正确？
\end{problem}

\vspace{0.8cm}

\begin{example}
    若 $(2-x)^{n}$ 的展开式中 $x^{2}$ 的系数为 240 ，则 $n=$
    \begin{tasks}(4)
        \task $4$
        \task $5$
        \task $6$
        \task $8$
    \end{tasks}
\end{example}

\subsection{两组乘法的二项式问题}

这一小节的题目也就是考试中常考的基础题。

\vspace{0.4cm}

\begin{example}
    $\left(x-\frac{1}{x}\right)(a+y)^{6}$ 的展开式中，含 $x^{-1} y^{4}$ 项的系数为 -15 ，则 $a=$
    \begin{tasks}(4)
        \task $-2$
        \task $-1$
        \task $\pm 1$
        \task $\pm 2$
    \end{tasks}
\end{example}

\begin{example}
    $\left(1+x^{2}\right)(1-x)^{5}$ 的展开式中 $x^{3}$ 的系数为 \underline{\hspace{2cm}} （用数字作答）。
\end{example}

\subsection*{练习}

\begin{enumerate}

    \item $\left(a-\frac{x}{y}\right)(x+y)^{7}$ 的展开式中 $x^{2} y^{5}$ 的系数为 56 ，则 $a$ 的值为 \underline{\hspace{2cm}} ．

    \item $\frac{y}{x}(x-2 y)^{5}$ 的展开式中 $x^{3} y^{2}$ 的系数为 \underline{\hspace{2cm}} （用数字作答）。

\end{enumerate}

\subsection{杨辉三角}

\begin{definition}
    杨辉三角是一个由二项式系数组成的三角形排列。它的第 $n$ 行的第 $k$ 个元素表示二项式系数 $C_n^k$，即：
    \begin{center}
        \begin{tabular}{*{11}{c}l}
              &   &   &   &    & 1 &    &   &   &   &   & $\quad (a+b)^0 = 1$                                   \\
              &   &   &   & 1  &   & 1  &   &   &   &   & $\quad (a+b)^1 = a + b$                               \\
              &   &   & 1 &    & 2 &    & 1 &   &   &   & $\quad (a+b)^2 = a^2 + 2ab + b^2$                     \\
              &   & 1 &   & 3  &   & 3  &   & 1 &   &   & $\quad (a+b)^3 = a^3 + 3a^2b + 3ab^2 + b^3$           \\
              & 1 &   & 4 &    & 6 &    & 4 &   & 1 &   & $\quad (a+b)^4 = a^4 + 4a^3b + 6a^2b^2 + 4ab^3 + b^4$ \\
            1 &   & 5 &   & 10 &   & 10 &   & 5 &   & 1 & $\quad (a+b)^5$                                       \\
        \end{tabular}
    \end{center}
\end{definition}

杨辉三角通常可以用来找找规律。特别是找出 $(a+b)^n$ 展开式中二项式系数最大的项是第几项，可以先手画几个找出规律。

\begin{example}
    在 $\left(x^{2}+\frac{3}{x}\right)^{n}\left(n \in \mathbb{N}^{*}\right)$ 的展开式中，二项式系数的最大值为 20 ，则系数的最大值为
    \begin{tasks}(4)
        \task $729$
        \task $1243$
        \task $1458$
        \task $2187$
    \end{tasks}
\end{example}

\begin{remark}
    二项式系数指的是 $C_n^k$ 的值，即杨辉三角上的值，而系数指的是这个式子展开后得到的每一项前的数。
\end{remark}

\begin{note}
    \vspace{3cm}
\end{note}

\begin{theorem}
    观察杨辉三角，可以发现，左右两边具有对称性。另外，每一行的和其实也有规律，分别是 $1, 2, 4, 8, 16, 32, ...$。
    \begin{enumerate}
        \item 二项式系数的对称性：$C_n^k = C_n^{n-k}$。
        \item 杨辉三角的每一行的和为 $2^n$，即 $\sum_{k=0}^{n} C_n^k = 2^n$。
        \item 奇数项和偶数项的二项式系数和均为 $2^{n-1}$。
    \end{enumerate}
\end{theorem}

\subsection{“三项式定理”}

\begin{example}
    多项式 $(x+y+z)^{4}$ 展开式中，$x y z^{2}$ 项的系数为
    \begin{tasks}(4)
        \task $12$
        \task $24$
        \task $48$
        \task $96$
    \end{tasks}
\end{example}

\begin{solution}
    只要当成取小球就好了，取 $x$ 小球 $1$ 个，还剩3个球，再取 $y$ 小球 $1$ 个，取 $z$ 小球 $2$ 次，即：$C_4^1 C_3^1 C_2^2.$
\end{solution}

\subsection*{练习}

\begin{enumerate}

    \item $(x+2 y-1)^{4}$ 的展开式中，$x^{2} y$ 的系数为
          \begin{tasks}(4)
              \task $24$
              \task $-24$
              \task $12$
              \task $-48$
          \end{tasks}

    \item 在 $(2+a+b)^{5}$ 的展开式中，$a^{2} b$ 的系数为
          \begin{tasks}(4)
              \task $30$
              \task $60$
              \task $90$
              \task $120$
          \end{tasks}

    \item 已知 $(1+x)^{n}$ 的展开式中第 3 项与第 7 项的二项式系数相等，则
          \begin{itemize}
              \item A. $n=8$
              \item B. $(1+x)^{n}$ 的展开式中 $x^{2}$ 项的系数为 56
              \item C. 奇数项的二项式系数和为 128
              \item D. $\left(1+x-y^{2}\right)^{n}$ 的展开式中常数项的系数为 1
          \end{itemize}
\end{enumerate}

\subsection{二项分布}

\begin{enumerate}

    \item 小刚参与一种答题游戏，需要解答 $A, B, C$ 三道题．已知他答对这三道题的概率分别为 $a, a, \frac{1}{2}$ ，且各题答对与否互不影响，若他恰好能答对两道题的概率为 $\frac{1}{4}$ ，则他三道题都答错的概率为
          \begin{tasks}(4)
              \task $\frac{1}{2}$
              \task $\frac{1}{3}$
              \task $\frac{1}{4}$
              \task $\frac{1}{6}$
          \end{tasks}


    \item 学生李明上学要经过 4 个路口，前三个路口遇到红灯的概率均为 $\frac{1}{2}$ ，第四个路口遇到红灯的概率为 $\frac{1}{3}$ ，设在各个路口是否遇到红灯互不影响，则李明从家到学校恰好遇到一次红灯的概率为
          \begin{tasks}(4)
              \task $\frac{7}{24}$
              \task $\frac{1}{4}$
              \task $\frac{1}{24}$
              \task $\frac{1}{8}$
          \end{tasks}
\end{enumerate}

\begin{enumerate}

    \item 某学校举行校园歌手大赛活动邀请了 6 位专家评委，在活动结束时邀请这 6 位专家站成一排合影留念，则下列说法正确的是
          \begin{tasks}(1) % Options are long, better as single column or (2)
              \task 若将专家甲、乙、丙三人从左到右按照身高递增的顺序排列，则共有 120 种排法
              \task 若专家甲、乙两人不相邻，则共有 480 种排法
              \task 若专家甲、乙、丙三人相邻，且甲在中间，则共有 72 种排法
              \task 若专家丙不在排头，丁不在排尾，则共有 480 种排法
          \end{tasks}

    \item 一个系统如图所示， $\mathrm{A}, B, C, D, E, F$ 为 6 个部件，其正常工作的概率都是 $\frac{1}{2}$ ，且是否正常工作是相互独立的，当 $\mathrm{A}, B$ 都正常工作或 $C$ 正常工作，或 $D$ 正常工作，或 $E$ ， $F$ 都正常工作时，系统就能正常工作，则系统正常工作的概率是 % Image: 2025_05_27_1339225c583797a01da5g-2
          \begin{tasks}(4)
              \task $\frac{55}{64}$
              \task $\frac{1}{64}$
              \task $\frac{1}{8}$
              \task $\frac{9}{64}$
          \end{tasks}

    \item 一场小型晚会有 3 个唱歌节目和 2 个相声节目，要求排出一个节目单．第一个节目和最后一个节目都是唱歌节目，有 \underline{\hspace{2cm}}种排法；前 3 个节目中要有相声节目，有 \underline{\hspace{2cm}}种排法。


    \item 总院计划安排小吴、小张等 8 名护士到 $A, B, C, D$ 四所医院实习，每两名护士到一所医院，考虑到个人意愿问题，小吴不安排 A 医院，小张不安排 $B$ 医院，则安排方法共有
          \begin{tasks}(4)
              \task $1200$ 种
              \task $1440$ 种
              \task $1600$ 种
              \task $1800$ 种
          \end{tasks}
\end{enumerate}

\subsection{方程的正整数解问题}

其实和之前提到的隔板法非常像，具体可以去看看那个小节。

\begin{example}
    方程 $x_{1}+x_{2}+x_{3}+x_{4}=7$ 的正整数解的个数为
    \begin{tasks}(4)
        \task $15$
        \task $35$
        \task $40$
        \task $20$
    \end{tasks}
\end{example}

\begin{solution}
    使用隔板法来解决正整数解分配问题。视为有七个球，现在要放三块板，隔开这些球，三块板隔出的四个区域分别代表 $x_1, x_2, x_3, x_4$ 的值，即总的情况数量有 $C_7^4$ 种。
\end{solution}

\vspace{2em}

\begin{example}
    已知 $x \in \mathbb{N}^{*}, y \in \mathbb{N}^{*}, z \in \mathbb{N}^{*}$ ，则关于 $x, y, z$ 的方程 $x+y+z=10$ 共有 \underline{\hspace{1.5cm}} 组不同的解．
    \begin{tasks}(4)
        \task $36$
        \task $45$
        \task $50$
        \task $24$
    \end{tasks}
\end{example}

\subsection*{练习}

\begin{enumerate}
    \item 四元一次方程 $x_{1}+x_{2}+x_{3}+x_{4}=10$ 的正整数解有 \underline{\hspace{2cm}}组，非负整数解有 \underline{\hspace{2cm}}组．

    \item 已知 $k$ 元一次方程 $x_{1}+x_{2}+x_{3}+\cdots+x_{k}=n\left(n \geq k, n \geq 2, n, k \in \mathbb{N}^{*}\right)$ 的正整数解的个数为 $\mathrm{C}_{n-1}^{k-1}$ ，则方程 $x_{1}+x_{2}+x_{3}+x_{4}=25$ 满足 $x_{i} \geq 2 i \quad(i=1,2,3,4)$ 的整数解的个数为
          \begin{tasks}(4)
              \task $35$
              \task $56$
              \task $84$
              \task $120$
          \end{tasks}
\end{enumerate}

\subsection{几何概型}

\begin{enumerate}
    \item 在等边 $\triangle A B C$ 中，$D$、$E$、$F$ 分别是 $A B$、$B C$、$C A$ 的中点．若在 $\triangle A B C$ 内随机取一点，则此点取自黑色部分的概率为 \underline{\hspace{2cm}}

    \item $\triangle A B C$ 是一个小型花园，阴影部分为一个圆形水池，且与 $\triangle A B C$ 的三边相切，知 $A B=10 \mathrm{~m}, A C=8 \mathrm{~m}, B C=6 \mathrm{~m}$ ．若从天空飘落下一片树叶恰好落入花园里，则落入水池的概率为 \underline{\hspace{2cm}} 。（ $\pi$ 取 3）

    \item 在区间 $(0,1)$ 随机取两个数，则两数之和小于 $\frac{3}{2}$ 的概率为
          \begin{tasks}(4)
              \task $\frac{1}{8}$
              \task $\frac{3}{8}$
              \task $\frac{5}{8}$
              \task $\frac{7}{8}$
          \end{tasks}
\end{enumerate}

\subsection{代值问题}

题目一般会给一个等式，左右两边是很明显的二项式展开的式子。而它会让你求 $\sum a_k$ 类似的值，这时候我们可以用赋值法处理（标答也是这样做的）。

\begin{example}
    若 $(x-2)^{5}=a_{0}+a_{1} x+a_{2} x^{2}+a_{3} x^{3}+a_{4} x^{4}+a_{5} x^{5}$ ，则 $a_{0}+a_{1}+a_{2}+a_{3}+a_{4}+a_{5}=$ \underline{\hspace{2cm}} ．

\end{example}

\begin{solution}
    观察到，代入 $x = 1$ 可以得到所有系数的和。
\end{solution}

\subsection*{练习}

\begin{enumerate}
    \item
          若 $(1-x)^{6}=a_{0}+a_{1}(x+1)+a_{2}(x+1)^{2}+\cdots+a_{6}(x+1)^{6}$ ，则 $a_{1}+a_{3}+a_{5}$ 的值为 \underline{\hspace{2cm}} ．

          \begin{remark}
              这道题还挺绕的，可以尝试一下赋值，让它等于 $1,-1,0$ 之类的试试。
          \end{remark}

          \begin{note}
              \vspace{2cm}
          \end{note}

    \item 有一种摸奖游戏，在一个口袋中有红球 4 个，黑球 2 个，小球除颜色外没有任何区别．
          规定：摸到红球记 1 分，摸到黑球记 -1 分．抽奖人从中一次摸 1 个球，摸 3 个球为一次抽奖，总分记为 $X$ ，若 $X \geq 0$ ，则获奖．
          \begin{enumerate}[(1)]
              \item 若每次取出后不放回，求获奖的概率；
              \item 若每次取出后放回，求 $X$ 的分布列和数学期望．
          \end{enumerate}
\end{enumerate}

\subsection{总练习}

\begin{enumerate}
    \item $\left(\sqrt{x}-\frac{2}{x}\right)^{6}$ 的二项展开式中的常数项为 \underline{\hspace{2cm}} ．

    \item 已知 $\left(\frac{2}{x}+\sqrt{x}\right)^{n}\left(n \in \mathbb{N}^{*}\right)$ 的展开式中第 7 项为常数项．
          \begin{enumerate}[(1)]
              \item 求 $n$ 的值；
              \item 从展开式中的所有项中任取两项，求取出的两项都是有理项的概率．
          \end{enumerate}
\end{enumerate}